\section{The Robustness--Specificity Tradeoff: Formal Analysis}

\subsection{Definitions}

Let $f_\phi$ be a pretrained text-to-image diffusion model and let
$\tilde{f}_\phi$ denote a concept-erased model (CEM) obtained by some
unlearning procedure applied to $f_\phi$. We work with the UNet noise
predictor $\epsilon_\theta$ and its associated text encoder
$\mathcal{Y}_\psi : \mathcal{T} \to \mathcal{P}$, where $\mathcal{T}$
is the text space and $\mathcal{P} \subseteq \mathbb{R}^d$ is the text
embedding space.

\begin{definition}[Concept activation region]
Let $\ell^*$ denote a fixed layer of the UNet (in our experiments,
\texttt{up\_blocks.1.attentions.1}) and let $t^* \in \{1,\ldots,T\}$
denote a representative denoising timestep. Define the
\emph{concept activation function}
\[
    \Phi_\theta : \mathcal{T} \longrightarrow \mathcal{Z}, \qquad
    \Phi_\theta(p) \;=\; \mathrm{A}_\theta(p,\, \ell^*,\, t^*),
\]
where $\mathrm{A}_\theta(p,\ell,t)$ denotes the cross-attention output
at layer $\ell$ and timestep $t$ under model $\theta$. The
\emph{activation region} of concept $c$ is
\[
    \mathcal{R}_c \;=\; \bigl\{\, \Phi_{\theta^*}(p) \;\big|\;
    p \in \mathcal{T},\; \mathbb{O}_\mathcal{X}(f_\phi(p), c) = 1
    \,\bigr\} \;\subseteq\; \mathcal{Z},
\]
where $\mathbb{O}_\mathcal{X} : \mathcal{X} \times \mathcal{C}
\to \{0,1\}$ is the ground-truth oracle that returns $1$ if concept
$c$ appears in image $x$.
\end{definition}

\begin{definition}[Semantic neighborhood]
The \emph{semantic neighborhood} of concept $c$ at radius $\delta > 0$
is
\[
    \mathcal{N}_\delta(c) \;=\; \bigl\{\, c' \in \mathcal{C}
    \;\big|\; c' \neq c, \;
    d_\mathcal{Z}\!\left(\mathcal{R}_c,\, \mathcal{R}_{c'}\right)
    < \delta \,\bigr\},
\]
where $d_\mathcal{Z}(A,B) = \inf_{a \in A, b \in B}\|a - b\|_2$
denotes the set distance in $\mathcal{Z}$.
\end{definition}

\begin{definition}[$\varepsilon$-robust unlearning]
\label{def:robust}
A CEM $\tilde{f}_\phi$ with unlearned parameters $\hat\theta$ achieves
\emph{$\varepsilon$-robust unlearning} of concept $c_u$ if, for all
prompts $p \in \mathcal{T}$ (including those that do not mention $c_u$
explicitly),
\[
    \Pr_{x \sim p_{\hat\theta}(\cdot \mid p)}\!\bigl[
    \mathbb{O}_\mathcal{X}(x, c_u) = 1\bigr] \;<\; \varepsilon.
\]
\end{definition}

\begin{definition}[$\gamma$-neighbor preservation]
\label{def:neighbor}
A CEM $\tilde{f}_\phi$ achieves \emph{$\gamma$-neighbor preservation}
with respect to $\mathcal{N}_\delta(c_u)$ if, for all
$c' \in \mathcal{N}_\delta(c_u)$,
\[
    \bigl\|\Phi_{\hat\theta}(p_{c'}) - \Phi_{\theta^*}(p_{c'})
    \bigr\|_2 \;<\; \gamma,
\]
where $p_{c'}$ is a canonical prompt for concept $c'$.
\end{definition}

\begin{definition}[Concept overlap]
The \emph{overlap coefficient} between concept $c_u$ and its
neighborhood $\mathcal{N}_\delta(c_u)$ is
\[
    \kappa(c_u, \delta) \;=\;
    \frac{
        \mathrm{vol}\!\left(
            \mathcal{R}_{c_u}^{(\rho)} \cap
            \bigcup_{c' \in \mathcal{N}_\delta(c_u)}
            \mathcal{R}_{c'}^{(\rho)}
        \right)
    }{
        \mathrm{vol}\!\left(\mathcal{R}_{c_u}^{(\rho)}\right)
    },
\]
where $\mathcal{R}_c^{(\rho)} = \{z \in \mathcal{Z} \mid
\inf_{r \in \mathcal{R}_c}\|z - r\|_2 \leq \rho\}$ denotes the
$\rho$-fattening of $\mathcal{R}_c$ for some margin $\rho > 0$.
\end{definition}

\subsection{Main Result}

\begin{assumption}[Lipschitz continuity of generation]
\label{asm:lip}
The mapping from activation space to concept generation probability
is $L$-Lipschitz: for all $z, z' \in \mathcal{Z}$,
\[
    \bigl|\Pr[\mathbb{O}_\mathcal{X}(f_\phi(z),c)=1]
    - \Pr[\mathbb{O}_\mathcal{X}(f_\phi(z'),c)=1]\bigr|
    \;\leq\; L\|z - z'\|_2.
\]
\end{assumption}

\begin{assumption}[Activation-editing constraint]
\label{asm:edit}
Any unlearning edit $\theta^* \mapsto \hat\theta$ that suppresses
$\mathcal{R}_{c_u}$ must displace activations by at least
$\Delta > 0$ within $\mathcal{R}_{c_u}$:
\[
    \forall\, p \;\text{s.t.}\;
    \Phi_{\theta^*}(p) \in \mathcal{R}_{c_u}:
    \quad
    \bigl\|\Phi_{\hat\theta}(p) - \Phi_{\theta^*}(p)\bigr\|_2
    \;\geq\; \Delta,
\]
where $\Delta \geq (1 - \varepsilon)/L$ is the minimum displacement
required to push generation probability below $\varepsilon$.
\end{assumption}

\begin{theorem}[Robustness--Specificity Tradeoff]
\label{thm:tradeoff}
Let Assumptions~\ref{asm:lip} and~\ref{asm:edit} hold. Suppose
$\tilde{f}_\phi$ achieves $\varepsilon$-robust unlearning of $c_u$
(Definition~\ref{def:robust}). Then for any neighbor concept
$c' \in \mathcal{N}_\delta(c_u)$ whose activation region satisfies
$\mathcal{R}_{c'}^{(\rho)} \cap \mathcal{R}_{c_u}^{(\rho)} \neq
\emptyset$, the neighbor preservation error is lower bounded by
\[
    \bigl\|\Phi_{\hat\theta}(p_{c'}) -
    \Phi_{\theta^*}(p_{c'})\bigr\|_2
    \;\geq\;
    \kappa(c_u,\delta) \cdot \Delta
    \;-\; \delta,
\]
and consequently, $\gamma$-neighbor preservation
(Definition~\ref{def:neighbor}) requires
\[
    \boxed{
    \gamma \;\geq\; \kappa(c_u,\delta)\cdot\frac{1-\varepsilon}{L}
    \;-\; \delta.
    }
\]
In particular, achieving both $\varepsilon$-robustness and
$\gamma$-neighbor preservation simultaneously requires
\[
    \varepsilon \;\geq\; 1 - \frac{L(\gamma + \delta)}{\kappa(c_u,\delta)}.
\]
\end{theorem}

\begin{proof}
We proceed in three steps.

\medskip
\noindent\textbf{Step 1: Required activation displacement for
$\varepsilon$-robustness.}

By Assumption~\ref{asm:lip}, for any prompt $p$ such that
$\Phi_{\theta^*}(p) \in \mathcal{R}_{c_u}$, the original model
generates concept $c_u$ with probability at least $1 - \varepsilon_0$
for some baseline $\varepsilon_0 < \varepsilon$. To reduce this
probability below $\varepsilon$ in the erased model, the activation
must be displaced by at least
\[
    \bigl\|\Phi_{\hat\theta}(p) - \Phi_{\theta^*}(p)\bigr\|_2
    \;\geq\; \frac{1 - \varepsilon}{L} \;=:\; \Delta,
\]
which is exactly the condition in Assumption~\ref{asm:edit}. This
establishes that $\varepsilon$-robustness imposes a uniform
displacement lower bound $\Delta$ on all activations within
$\mathcal{R}_{c_u}$.

\medskip
\noindent\textbf{Step 2: Displacement propagates to overlapping
neighbor regions.}

Let $c' \in \mathcal{N}_\delta(c_u)$ and let
$\mathcal{O} = \mathcal{R}_{c_u}^{(\rho)} \cap
\mathcal{R}_{c'}^{(\rho)}$ denote the overlap region. By hypothesis,
$\mathcal{O} \neq \emptyset$.

Since the unlearning edit must displace all activations in
$\mathcal{R}_{c_u}^{(\rho)}$ by at least $\Delta$ (from Step 1), and
since $p_{c'}$ is chosen such that
$\Phi_{\theta^*}(p_{c'}) \in \mathcal{R}_{c'}$, continuity of
$\Phi_{\hat\theta}$ implies that the displacement at $p_{c'}$ is
at least
\[
    \bigl\|\Phi_{\hat\theta}(p_{c'}) - \Phi_{\theta^*}(p_{c'})
    \bigr\|_2
    \;\geq\;
    \frac{\mathrm{vol}(\mathcal{O})}{\mathrm{vol}(\mathcal{R}_{c_u}^{(\rho)})}
    \cdot \Delta \;-\; d_\mathcal{Z}(\mathcal{R}_{c_u}, \mathcal{R}_{c'}).
\]
By Definition 2, $d_\mathcal{Z}(\mathcal{R}_{c_u},\mathcal{R}_{c'})
< \delta$ for all $c' \in \mathcal{N}_\delta(c_u)$. Substituting the
definition of $\kappa(c_u,\delta)$ yields
\[
    \bigl\|\Phi_{\hat\theta}(p_{c'}) - \Phi_{\theta^*}(p_{c'})
    \bigr\|_2
    \;\geq\;
    \kappa(c_u,\delta) \cdot \Delta \;-\; \delta.
\]

\medskip
\noindent\textbf{Step 3: Deriving the tradeoff inequality.}

Substituting $\Delta = (1-\varepsilon)/L$ from Step 1:
\[
    \bigl\|\Phi_{\hat\theta}(p_{c'}) - \Phi_{\theta^*}(p_{c'})
    \bigr\|_2
    \;\geq\;
    \kappa(c_u,\delta) \cdot \frac{1-\varepsilon}{L} \;-\; \delta.
\]
For $\gamma$-neighbor preservation (Definition~\ref{def:neighbor})
to hold, we need the left-hand side to be less than $\gamma$, which
requires
\[
    \gamma \;\geq\;
    \kappa(c_u,\delta) \cdot \frac{1-\varepsilon}{L} \;-\; \delta.
\]
Rearranging for $\varepsilon$:
\[
    \kappa(c_u,\delta) \cdot \frac{1-\varepsilon}{L}
    \;\leq\; \gamma + \delta
    \quad\Longrightarrow\quad
    1 - \varepsilon \;\leq\;
    \frac{L(\gamma+\delta)}{\kappa(c_u,\delta)}
    \quad\Longrightarrow\quad
    \varepsilon \;\geq\;
    1 - \frac{L(\gamma+\delta)}{\kappa(c_u,\delta)}.
\]
This completes the proof.
\end{proof}

\begin{corollary}[Impossibility of simultaneous perfect erasure
and perfect preservation]
\label{cor:impossibility}
If $\kappa(c_u, \delta) > 0$ and $\delta < \kappa(c_u,\delta)/L$,
then there is no unlearning method that simultaneously achieves
$\varepsilon = 0$ (perfect erasure) and $\gamma = 0$ (perfect
neighbor preservation).
\end{corollary}

\begin{proof}
Setting $\varepsilon = 0$ in Theorem~\ref{thm:tradeoff} gives
$\gamma \geq \kappa(c_u,\delta)/L - \delta > 0$ under the stated
conditions. Hence $\gamma = 0$ is infeasible.
\end{proof}

\begin{remark}[Connection to STEREO's robustness mechanism]
Theorem~\ref{thm:tradeoff} provides a formal explanation for why
STEREO achieves robustness despite using only $K=2$ adversarial
prompts in the STE stage. The learned token embeddings
$v_1^*, v_2^*$ (Eq.~4 in \cite{stereo}) are not natural-language
prompts but optimized vectors in $\mathcal{P}$ that lie near the
hardest-to-suppress directions of $\mathcal{R}_{c_u}$. The REO
compositional objective then suppresses the entire subspace spanned
by these directions, effectively covering $\mathcal{R}_{c_u}$ with
only two reference points. This broad suppression corresponds to
a large $\Delta$ in Assumption~\ref{asm:edit}, which by
Theorem~\ref{thm:tradeoff} directly implies large collateral
displacement in $\mathcal{R}_{c'}$ for all
$c' \in \mathcal{N}_\delta(c_u)$ with nonzero overlap
$\kappa(c_u,\delta) > 0$. Our empirical observation that deer and
zebra generation degrades after horse erasure (Section~3.2) is
consistent with $\kappa(\texttt{horse}, \delta) > 0$ for these
neighboring animal concepts.
\end{remark}

\begin{remark}[Tightness]
The bound in Theorem~\ref{thm:tradeoff} is tight when the overlap
region $\mathcal{O}$ is exactly $\kappa(c_u,\delta)$ fraction of
$\mathcal{R}_{c_u}^{(\rho)}$ and the displacement field induced by
the unlearning edit is uniform over $\mathcal{R}_{c_u}^{(\rho)}$.
In practice, the displacement may be non-uniform (as suggested by
the layer-wise cross-attention similarity results in
Figure~\ref{fig:attn_similarity}), in which case the bound is
conservative. Tightening it by characterizing the displacement
geometry more precisely is left for future work.
\end{remark}